\section{Analysis}\label{sec:analysis}

We analyse the final DeBERTa 10-seed ensemble (argmax) and compare it with the paired
Qwen runs of \S\ref{sec:llm}, all on dev.

\paragraph{The model, not label noise, fails.}
Of the ensemble's 308 predictions, 137 (44.5\%) fell outside the item's reference set, and
53 of these were on items where all three annotators agreed. Annotator disagreement is real
and informative \citep{uma2021learning}, but it does not explain the gap. Scored the way
the test set is scored, each annotator against the other two reaches 0.643--0.766 on \Stwo{}
(mean 0.684), while the ensemble reaches 0.380--0.381.

\paragraph{Commitment, not only topic.}
The largest error pairs (reference $\rightarrow$ prediction) are Implicit $\rightarrow$
Explicit (25), General $\rightarrow$ Implicit (20) and General $\rightarrow$ Explicit (20).
All three concern how committed an answer is, more than whether it stays on topic. Length
matters too. The 20 inputs that are still cut at 1,024 tokens land in-set only 20\% of the
time (0.580 for the rest), and the shorter half of dev scores 0.400 on \Stwo{} against 0.283
for the longer half.

\FigureExample
A typical error shows how much an answer can depend on world knowledge
(Figure~\ref{fig:example}). Asked whether he would campaign against Senator Lieberman, the
speaker replied \emph{``I'm going to stay out of Connecticut.''} Seeing this as a near-answer
requires knowing that Lieberman is a senator from Connecticut, which is our reading of the
item rather than something we measured, and errors of this kind are what led us to the
knowledge hypothesis. Qwen ranks Dodging first, but logit adjustment then picks the rarer
Declining to answer. A prior correction that raises macro-F1 over the whole set can still
cost an individual item.

\paragraph{Where the 8B classifier gains.}
Table~\ref{tab:analysis-classes} shows that our registered prediction was wrong. Per-class
F1 rises in 7 of the 9 classes, most in Claims ignorance (+0.403), Declining (+0.220),
Dodging (+0.103) and Deflection (+0.089). The commitment classes move least (General +0.059,
Explicit +0.057, Implicit +0.040), and neither model ever predicts Partial/half-answer. The
in-set rate rises on unanimous items (+0.076) about as much as on contested ones (+0.105 and
+0.048). The biggest gains are on the Non-Reply classes, which depend on recognising what
kind of statement an answer is, while General stays the weakest frequent class (F1 0.322).

\begin{table}[t]
\centering
\small
\begin{tabular}{@{}lccc@{}}
\toprule
 & DeBERTa & Qwen & $\Delta$ \\
\midrule
\multicolumn{4}{@{}l}{\emph{per-class F1}} \\
Explicit            & 0.672 & 0.729 & +0.057 \\
Implicit            & 0.403 & 0.443 & +0.040 \\
Dodging             & 0.484 & 0.587 & +0.103 \\
General             & 0.263 & 0.322 & +0.059 \\
Deflection          & 0.287 & 0.377 & +0.089 \\
Partial/half-answer & 0.000 & 0.000 & +0.000 \\
Declining to answer & 0.381 & 0.601 & +0.220 \\
Claims ignorance    & 0.296 & 0.700 & +0.403 \\
Clarification       & 0.667 & 0.524 & $-$0.144 \\
\midrule
\multicolumn{4}{@{}l}{\emph{in-set rate by distinct annotator labels}} \\
one ($n=125$)   & 0.529 & 0.605 & +0.076 \\
two ($n=150$)   & 0.493 & 0.598 & +0.105 \\
three ($n=33$)  & 0.688 & 0.736 & +0.048 \\
\bottomrule
\end{tabular}
\caption{Where the backbone gain falls, for single models averaged over the same 10 seeds
(dev \Stwo{}). The registered prediction, a gain concentrated on the commitment boundary and
on agreed items, does not hold.}
\label{tab:analysis-classes}
\end{table}

\paragraph{Ranking and the decision rule.}
Both models rank better than they decide. For Qwen (seeds 0--2), an acceptable label is the
top prediction for 60.4\% of items and in the top three for 93.1\%, and the in-set rate rises
from 0.40 to 0.87 across fifths of dev from least to most confident, so most remaining errors
are a choice among a few candidates the model already ranks highly. Raw Qwen leans towards frequent classes. On seeds 0--2 it predicted General
for only 24 items, although General is in 113 reference sets. Logit adjustment corrects
this, adding +0.062 per Qwen model (0.476 to 0.538) against +0.005 for DeBERTa. Mixing Qwen
and DeBERTa probabilities, with the weight fitted inside the nested CV, adds nothing (0.543
against 0.549 for Qwen alone).
